% Generated draft: requires mathematical and source review.
\documentclass[11pt]{article}
\usepackage[margin=25mm]{geometry}
\usepackage{fontspec}
\setmainfont{Latin Modern Roman}
\usepackage{amsmath,amssymb,mathtools,unicode-math}
\setmathfont{Latin Modern Math}
\usepackage{enumitem,needspace,xurl,hyperref}
\setlength{\parindent}{0pt}
\setlength{\parskip}{5pt}
\setlength{\emergencystretch}{3em}
\newcommand{\N}{\mathbb{N}}
\newcommand{\Z}{\mathbb{Z}}
\newcommand{\Q}{\mathbb{Q}}
\newcommand{\R}{\mathbb{R}}
\newcommand{\C}{\mathbb{C}}
\newcommand{\F}{\mathbb{F}}
\newcommand{\A}{\mathbb{A}}
\newcommand{\PP}{\mathbb{P}}
\newcommand{\E}{\mathbb{E}}
\newcommand{\eps}{\varepsilon}
\newcommand{\dd}{\,\mathrm{d}}
\begin{document}
\section*{UnsolvedMath ID: 701354}
\textbf{Source problem code:} catalogue entry
\subsection{Definitions and mathematical statement}
Let \(p\) be a prime, let \(F\) be a finite field of characteristic \(p\), and let \(R\) be a finite-dimensional commutative associative \(F\)-algebra, with no multiplicative identity required. Assume that \(R\) is nilpotent: there exists an integer \(m\geq 1\) such that every product of \(m\) elements of \(R\) is zero. Define
\[
R^{(p)}=\{r^p:r\in R\}.
\]
This set is an \(F\)-subalgebra: in characteristic \(p\), commutativity gives \((r+s)^p=r^p+s^p\) and \((rs)^p=r^ps^p\), and every scalar in the finite field \(F\) has a \(p\)-th root in \(F\).

The question attributed in the supplied statement to N. Eggert is whether, for every such \(F\) and \(R\),
\[
\dim_F R\ \geq\ p\,\dim_F R^{(p)}.
\]
The supplied record further states that an affirmative answer would solve Kourovka Notebook Problem 11.44 in the case where its factors \(A\) and \(B\) are abelian. The statement of Problem 11.44 and the meaning of those factors are not supplied, so that claimed implication is retained as source context without a formal reconstruction.
\subsection{Short English statement}
For every finite-dimensional commutative associative nilpotent algebra over a finite field of characteristic \(p\), must its dimension be at least \(p\) times the dimension of its subalgebra of \(p\)-th powers? The record attributes the question to N. Eggert and lists L. S. Kazarin as proposer. Its \emph{open} status is supplied registry metadata, not an independently verified assessment.
\subsection{Sources}
\begin{itemize}
\item Local UnsolvedMath registry, supplied complete source record: permanent ID \(701354\), exact source problem code \texttt{catalogue entry}, title \emph{Kourovka Notebook Problem 17.52}. Upstream revision: \texttt{d371ee917ac9647fbb06cd6fc529deec6f3dda82}. The record was created and updated on \texttt{2026-10-16T14:00:00Z} and labels the problem open; this is supplied metadata, not independent verification.
\item \emph{The Kourovka Notebook (Unsolved Problems in Group Theory)}, Problem 17.52, 17th issue (2010); the supplied background also identifies the 21st issue edition (2026), \texttt{arXiv:1401.0300v48}. The collection metadata names E. I. Khukhro and V. D. Mazurov as editors.
\item Attribution as supplied: L. S. Kazarin is listed as proposer, with proposed year 2010; the mathematical statement begins with the attribution ``N. Eggert.''
\item The supplied statement refers to Kourovka Notebook Problem 11.44 and claims that an affirmative answer resolves its case with abelian factors \(A\) and \(B\). That problem's statement is not included in the supplied evidence.
\end{itemize}
\end{document}
